\section{Conclusion}

\subsection{What This Document Specifies}

\EM{} is specified here as a mobile application that treats a photograph as a
latent image awaiting development rather than as a finished picture awaiting a
filter. The specification is complete enough to build from: the imaging model is
given in closed form with numerical defaults and validation tolerances; the GPU
architecture is given with the fusion strategy, precision policy, and memory
budget that make it meet its latency targets; the application architecture,
interface, and persistence schema are given at implementation detail; and the
programme is given with phases, staffing, cost, and a risk register that names
owners and triggers.

The technical argument reduces to a single structural claim. A lookup table is a
pointwise function, and the phenomena that identify photographic film ---
halation, grain, interlayer inhibition, shoulder compression of scene highlights
--- are not pointwise. They are spatial, stochastic, or dependent on radiometric
information that a display-referred image no longer contains. No amount of
refinement to a pointwise operator on a rendered image closes that gap, because
the gap is one of representational class rather than of quality.

Carrying scene-referred data through a chain of physically motivated transforms
does close it, and the performance analysis in
Section~\ref{sec:performance} shows that doing so costs 2\,ms of GPU time per
preview frame on current mobile silicon. The obstacle was never computational.

\subsection{What Would Falsify It}

A design document that cannot be wrong is not a design document. Three concrete
outcomes would invalidate the approach as specified:

\begin{enumerate}
\item \textbf{The measurements do not close.} If AC-1 through AC-4 cannot be
met for a majority of target stocks with available or obtainable reference data,
then the model is being fitted to guesses and its claim to physical grounding is
rhetorical. Risk T-01, checked at month 4.
\item \textbf{Correctness does not produce conviction.} If the perceptual panel
continues to identify the renders reliably while every numerical acceptance
criterion passes, the model is measuring the wrong things. Risk T-02, checked at
month 7 and again at month 15.
\item \textbf{The distinction does not matter commercially.} If non-expert users
show no preference between a physically simulated render and a strong LUT
baseline, the product is correctly built and incorrectly aimed. Risk P-01,
checked during Phase 1.
\end{enumerate}

Each has a defined trigger and a prepared response. That is the difference
between a plan and an intention.

\subsection{The Argument in One Sentence}

Every quantity in this document that a user can adjust corresponds to something
that exists in a photographic laboratory, and every equation that consumes those
quantities can be checked against a densitometer. If both remain true through
implementation, the product's claim is not a marketing position but a
verifiable property --- and a photograph made with it will not have been
filtered. It will have been developed.
